\begin{center}
	\chapter{CONCLUSIONS AND FUTURE WORK}
\end{center}
\justifying

\section{Conclusions}

This work addressed the allocation of computation in linear recurrent layers
whose state-transition operator is a product of generalized Householder factors,
where the number of factors per token is conventionally a global hyperparameter.

The principal result is an order-requirement model for such layers. The number of
factors a symbol requires is determined by the minimum, over all faithful
representations of the group generated by the entire alphabet, of the rank of the
symbol's deviation from the identity. This quantity is not an attribute of the
symbol: a five-cycle requires four factors when the alphabet generates the
symmetric group $S_5$ and two when it generates the alternating group $A_5$,
because the latter admits a three-dimensional faithful representation as the
rotation group of the icosahedron. The model was verified by constructing that
representation from geometry and confirming every property it depends upon, with
a worst reconstruction error of $1.31 \times 10^{-15}$. Its central prediction,
that the generated group rather than the alphabet size governs the required
order, was confirmed experimentally: at a fixed factor count of three, alphabets
generating $A_5$ reach token accuracies between $0.9952$ and $0.9994$ while
alphabets generating $S_5$ reach $0.1094$ and $0.0307$, with the largest alphabet
tested falling in the former group and one of the smallest in the latter.

Two reachability bounds were established by explicit construction over all $120$
elements of $S_5$ at every factor count. The rank bound requires the factor count
to be at least the transposition length of the target. The determinant-parity
bound shows that when the factor strength is held at its reflection value, only
factor counts congruent to the transposition length modulo two admit an exact
representation, so that half of all targets are unreachable at any given count.
This second bound is not a peripheral observation: it determines that a mechanism
varying the effective order must scale a learnable factor strength toward zero
rather than select among fixed values, and the architecture proposed here is
constructed accordingly.

The practical significance of the order-requirement model is that it renders the
factor count a hyperparameter that cannot be selected reliably by inspection.
Setting it correctly requires the minimum faithful representation dimension of
the group generated by the data, which changes discontinuously as the alphabet
changes and is not observable from a dataset. A layer that infers its own order
therefore removes a hyperparameter that has no reliable manual setting, and this
argument is independent of any reduction in computation.

On the mechanism itself, an adaptive-order layer was implemented in which a
monotone halting gate assigns a per-token factor count and is trained under a
penalty on expected computation with no supervision on the allocation. Under an
oracle allocation supplied from ground truth, the layer attains the accuracy of
the cheapest adequate fixed order at $1.150$, $1.300$ and $1.752$ factors per
token where that fixed order costs four, corresponding to reductions in required
work of $3.48\times$, $3.08\times$ and $2.29\times$. Trained end to end across
three random seeds, the halting gate reproduced this allocation exactly in one
seed, assigning one factor to every requirement-one token and four to every
requirement-four token across $256{,}000$ evaluated tokens, at two independent
budget weights. The allocation was recovered rather than approximated.

A supporting result concerns measurement. Across fifty training runs, outcomes on
the harder alphabets were found to be bimodal rather than concentrated, with
escape from the loss plateau occurring in between $10\%$ and $30\%$ of runs and
selected by kernel non-determinism alone. The appropriate summary statistic on
such configurations is the escape fraction over repeated runs, and approximately
twenty runs per condition are required to distinguish escape rates of practical
interest. Setting these rates against the reduction available on each
configuration establishes that the configurations offering reductions above
$2.29\times$ train reliably, while those that train unreliably offer $1.28\times$
or less, so that the instability of the base architecture and the operating
region of the proposed mechanism are disjoint.

Two limitations bound the scope of these conclusions. All reductions reported
here are in required work per token, measured under a dense execution strategy in
which disabled factors are materialised and multiplied by zero; the adaptive and
full-cost conditions differ in elapsed time by less than one percent, and no
reduction in elapsed time is claimed. And the end-to-end allocation result rests
on three seeds at a single operating point on a single task family.

\section{Future Work}

\subsection{Reliability of the learned allocation}

The immediate objective is to establish the conditions under which the halting
gate converges to the correct allocation reliably. The evidence points to the
gate's initialisation rather than to the strength of the budget penalty: at the
largest penalty weight tested, the non-converging seeds held cost at exactly the
full factor count rather than at an intermediate value, which indicates a gate
that is not responding to the penalty rather than one responding insufficiently.
Two properties of the initialisation are relevant. The output weight of the gate
network is initialised at zero, which makes the gate constant with respect to its
input until that weight moves, so that the gate can raise or lower one global
order but cannot allocate. And the output bias places the halting nonlinearity at
a point where its derivative is small, attenuating the gradient that would move
that weight. Both are addressed by parameters now exposed for controlled study:
the initialisation bias of the gate, and a learning rate applied to the gate
independently of the remainder of the model. These will be evaluated one at a
time over ten seeds per condition, under the measurement protocol of Section 3.7.

\subsection{Inference of the order requirement from data}

A sharper form of the order-requirement model can be tested once the allocation
is reliable. The requirement for a five-cycle differs between an $A_5$ alphabet
and an $S_5$ alphabet, yet the symbol itself is identical in both. A gate that
assigns two factors under the first and four under the second would be inferring
a global algebraic property of the alphabet from purely local statistics, which
constitutes a considerably stronger result than a correlation between the
allocation and a measure of difficulty.

\subsection{Ragged execution and elapsed-time measurement}

Converting the measured reduction in required work into a reduction in elapsed
time requires an execution strategy in which disabled factors are not
materialised. Implementing this, and benchmarking it both at batch size one and
at the batch sizes typical of training, is the work that would substantiate the
efficiency argument. The benchmark is to be conducted in a manner that admits the
outcome that batching absorbs the saving, which is a possible and informative
result.

\subsection{Generality across groups and tasks}

The order-requirement model is stated for arbitrary finite groups, but has been
verified empirically within $S_5$ and $A_5$. Extending the empirical
verification to $S_3$, $S_4$ and $A_4$ would confirm that the model predicts the
required factor count across groups rather than within a single one. Beyond the
permutation setting, evaluation on a state-tracking task not designed for this
study, such as tracking board configurations in a two-player game, would
establish whether the mechanism transfers to tasks whose state structure is not
a group word problem by construction.

\subsection{Consolidation of the experimental record}

Several results reported here rest on a single training run per condition,
specifically the closure grid of Section 4.4, the matched-computation sweep of
Section 4.5, and the error characterisation of Section 4.7. Repetition of these
under the protocol of Section 3.7, with the number of runs reported alongside
each figure, is scheduled before the final evaluation. A broader survey of the
linear-attention and state-space literature is to be completed over the same
period.
\null
